\chapter{Wykorzystane technologie}
\section{Wprowadzenie}
W tym rozdziale omówione zostaną technologie (języki programowania, protokoły, usługi) wykorzystane do realizacji projektu.

\section{Język programowania Go}
\subsection{Geneza oraz użycie Go}
Utworzony w 2006 roku przez trzech pracowników Google - Roba Pike'a, Kena Thompsona oraz Roberta Griesemera - Go to język programowania charakteryzujący się prostotą składni, która przypomina język C, z wydajnością języków kompilowanych, takich jak C++ lub Rust. Go jest obecnie popularne w projektach backendowych oraz chmurowych, wykorzystywane przez przykładowo Google, Dropbox, Uber, PayPal, Soundcloud, BBC, Docker, Kubernetes, American Express oraz Netflix. 

\subsection{Cechy Go}
Go jest językiem statycznie typowanym, o typowaniu silnym - gdzie typ zmiennej (przykładowo, liczba całkowita, liczba zmiennoprzecinkowa) znany jest przed uruchomieniem programu, a każda zmienna ma odgórnie ustalony typ w kodzie. 

Jest również językiem kompilowanym - kod nie jest na bieżąco interpretowany, jak w przypadku przykładowo Pythona, dzięki czemu większość błędów jest wykrywana zanim program zostanie uruchomiony. Zamiast tego tworzony jest w pełni przenośny, uruchamialny plik binarny, który uruchomiony może być na systemach Windows, Linux lub macOS. 

Go wyposażone jest w prostą, jednak wysoko efektywną bibliotekę standardową, co pozwala na minimalizację zewnętrznych zależności. Słynie również ze swojego modelu współbieżności, opartego na \enquote{goroutines} w przeciwieństwie do tradycyjnie wykorzystywanych wątków. Go wymaga jawnej obsługi błędów w przypadku każdej funkcji, która może zwrócić błąd - co ułatwia gwarancję poprawnego funkcjonowania budowanej aplikacji.

\subsection{Model współbieżności Go}
Go wspiera współbieżne wykonywanie zadań w ramach jednego programu, dzięki swojej alternatywie dla wątków - goroutines (dalej: gorutyny). Język wyposażony jest w swój własny program planujący, który przypisuje osobne \enquote{gorutyny} wątkom systemu operacyjnego - co pozwala na przypisanie kilku pojedynczemu.

Typowy wątek wymaga stosu o wielkości jednego do dwóch megabajtów, w czasie gdy gorutyna potrzebuje zaledwie dwóch kilobajtów. Różnicą jest również to, w jaki sposób dochodzi do dzielenia się danymi pomiędzy wątkami. Tradycyjne wątki odczytują i nadpisują współdzieloną pamięć, w czasie gdy Go zachęca do używania tzw. kanałów, przez które przesyłane mogą być dane lub sygnały, ułatwiając omijanie zjawiska wyścigu, to jest błędów wywołanych pseudolosową kolejnością wykonywanych współbieżnych zadań.

Go posiada również funkcjonalność zwana kontekstem, pozwalająca na wyznaczanie terminów oraz wysyłanie żądań między skoordynowanymi elementami programów.

\subsection{Uzasadnienie wyboru Go do wykonania projektu}
Go zostało wybrane przedewszystkim z uwagi na jego znajomość przez autora pracy - ale również prostotę składni i pracy z tym językiem. Kod jest przejrzysty i prosty w podziale na osobne elementy, a dzięki wymaganej kompilacji programu przy każdej próbie jego uruchomienia wykrywanie większości błędów odbywało się samoistnie. Biorąc pod uwagę założenia projektu, język umożliwił prostotę wdrożenia - plik końcowy uruchamialny jest na systemach Windows, Linux i macOS. 

Dzięki swojej bibliotece standardowej, cały projekt korzysta zaledwie z trzech zewnętrznych zależności:

\begin{small}
\begin{longtable}{>{\raggedright\arraybackslash}p{4.5cm} >{\raggedright\arraybackslash}p{3cm} >{\raggedright\arraybackslash}p{6cm}}
\caption{Biblioteki zewnętrzne wykorzystane w projekcie}
\label{tab:external_libraries} \\
\toprule
\textbf{Biblioteka} & \textbf{Zastosowanie} & \textbf{Rola w projekcie} \\
\midrule
\endfirsthead

\toprule
\textbf{Biblioteka} & \textbf{Zastosowanie} & \textbf{Rola w projekcie} \\
\midrule
\endhead

\bottomrule
\endfoot

github.com/goccy/go-yaml & konfiguracja & Odczyt i przetwarzanie pliku konfiguracyjnego \texttt{config.yaml} w formacie YAML. \\
github.com/terminalstatic/go-xsd-validate & walidacja XML & Walidacja wygenerowanych faktur XML wobec oficjalnego schematu XSD faktur ustrukturyzowanych FA(3). \\
modernc.org/sqlite & baza danych & Sterownik do bazy danych SQLite. \\
\end{longtable}
\end{small}

Go, podobnie jak C, posiada wskaźniki pamięci, które wykorzystane zostały w projekcie do opcjonalnych pól faktur oraz odróżniania brakujących danych od danych, które są zwyczajnie puste. Najważniejszym elementem jednak są wspomniane w poprzedniej subsekcji gorutyny, umożliwiające aplikacji współbieżną obsługę żądań HTTP oraz cykliczne skanowanie bazy danych w poszukiwaniu rekordów do przetworzenia. Sama wysyłka faktur oraz webhooków również jest współbieżna, a limit maksymalnych asynchronicznych zadań narzucany jest przez pojemność kanałów. Aplikacja oczekuje zakończenia obecnej próby przed rozpoczęciem kolejnego cyklu działania poprzez użycie \texttt{WaitGroup}, a użycie kontekstu pozwoliło na wprowadzenie mechanizmu bezpiecznego zamykania programu.

\section{Baza danych SQLite}

\section{Protokół HTTP oraz styl architektoniczny REST API}

\section{Format JSON}

\section{Format XML oraz schemat XSD}

\section{HTMX}

\section{Tailwind CSS}

\section{Docker}

\section{Środowisko Krajowego Systemu e-Faktur}
